\chapter{Relation to the industry design}
\label{ch:industry}

Public material from a 2026 TM~Forum Moonshot Catalyst on business-aware GNN healing
(C26.0.965) describes a production-grade system of the same shape as this repository. We built
ours independently from the blueprint, then compared. The comparison is useful in both directions:
it shows which parts of the industry system are mathematics that anyone can reproduce, and it
shows what we deliberately left out.

\section{Stage by stage}
\begin{center}\small
\begin{tabular}{@{}p{2.2cm}p{5.8cm}p{6.2cm}@{}}
\toprule
Stage & Industry design & This repository \\
\midrule
Trigger & Service-layer $z$-score against a clean multi-week baseline with a confirmation gate; the GNN does not detect. &
Same service-layer trigger implemented as a baseline (\code{service\_zscore\_trigger}); the GNN degradation field \eqref{eq:detection} is offered as a graph-aware alternative and measured against it. \\[2pt]
Localise & Link prediction: anomaly nodes live in the property graph and the model scores ``originates-from'' edges, traversed in reverse from symptoms to origin. &
Node scoring with a softmax over all nodes \eqref{eq:rca}. Equivalent for a single incident (Remark~\ref{rem:linkpred}); transposed relations in \eqref{eq:rgcn} are the reverse traversal. \\[2pt]
Prioritise & Utility $=$ profile weight $\times$ revenue-density index $\times$ service impact, banded into a priority; a change-risk concurrency policy executes one major change at a time. &
$R(v)=\sum_u w_u\,\mathrm{rev}_u\,\hat d_u$ over the blast radius, divided by resource cost, under a budget with one change per domain \eqref{eq:selection}. Their utility is the single-node special case. \\[2pt]
Execute & TMF921 v5 healing intent decomposed into TMF641 v5 service orders, executed by a vendor orchestration engine. &
TMF921- and TMF641-shaped objects (field names, not conformance) executed by a mock orchestrator against the simulator. \\[2pt]
Verify & A reflection stage re-runs inference, reports complies/degrades, re-triggers without human escalation. &
Re-simulation with the remediation applied; if the field does not recover, re-trigger with the next ranked candidate \eqref{eq:exposure}. \\[2pt]
Training data & A 13-stage schema-faithful generator: scripted temporal arcs, domain-specific signal fingerprints, topology-driven blast radii, root-cause vs symptom labels. &
A diffusion model on the causal graph \eqref{eq:propagation} with type-specific effect vectors \eqref{eq:observation}; origin and latent field as labels. Less schema fidelity, more physics. \\[2pt]
Scale substrate & A managed temporal graph database and a managed GNN training service. &
NVIDIA path: WholeGraph, cuGraph-PyG, cuDF (Chapter~\ref{ch:nvidia}); durable store left to the reader. \\
\bottomrule
\end{tabular}
\end{center}

\section{What transfers and what does not}
Everything in the Trigger, Localise, Prioritise and Verify rows is mathematics and is in this
repository with tests. The Execute row is plumbing that any orchestrator can provide. The two
things the industry system has that this repository does not are a production-identical schema for
its synthetic corpus, which makes fine-tuning on real incidents a no-translation step, and the
standards contributions themselves (an intent-ontology extension, a new assessment questionnaire,
a proposed lead-time indicator). Neither affects the claim of Chapter~\ref{ch:claim}; both matter
for deployment.

\section{One honest difference in the detection story}
The industry trigger is a service-layer statistic. That is a sensible production choice: it is
cheap, explainable, and needs no labels. Our results (Chapter~\ref{ch:experiments}) test whether a
graph-aware detector fires earlier on the same episodes with the same calibrated false-alarm
budget, because it reads the \emph{pattern across neighbours} rather than a single node's
deviation. That is an extension of the industry design, not a replication of it, and we report it
as such, including the cases where it does not win.
